\chapter{团队介绍}

\noindent\textit{\small 注：团队成员姓名为示例占位，正式版本请替换为实际成员信息。}

\section{团队构成与分工}

本项目团队为\hlgreen{在校本科生 / 研究生创业团队}，成员分别来自计算机、农业植保与市场营销相关专业，形成「技术开发—农技内容—落地推广」的完整分工。团队成员均为在校学生，具备充足的时间投入与学校创新创业平台、导师实验室的资源支持。

\begin{table}[H]
  \centering
  \small
  \renewcommand{\arraystretch}{1.4}
  \caption{核心团队成员（学生）}
  \begin{tabular}{p{2.6cm}p{3.4cm}p{8.2cm}}
    \toprule
    角色 & 姓名（示例） & 专业背景与职责 \\
    \midrule
    项目负责人 & 张 伟 & 农业信息化相关专业本科生，负责项目整体统筹、商业方案设计与对外沟通 \\
    算法负责人 & 李 明 & 计算机科学与技术专业本科生，负责图像识别模型训练、轻量化压缩与端侧部署 \\
    产品与开发 & 王 芳 & 软件工程专业本科生，负责小程序与 App 开发、交互设计与用户增长 \\
    农技内容 & 赵 强 & 植物保护专业本科生，负责病虫害知识库整理、防治方案撰写与专家对接 \\
    市场运营 & 孙 婷 & 市场营销专业本科生，负责县域推广、农户调研与合作渠道拓展 \\
    \bottomrule
  \end{tabular}
\end{table}

\section{指导教师与专家顾问}

\begin{itemize}[leftmargin=1.6em]
  \item \hlgreen{第一指导教师（植保方向）}：某农业大学植保学院教师，长期从事主要农作物病虫害监测预警研究，负责知识库与防治方案审校。
  \item \hlgreen{第二指导教师（人工智能方向）}：某高校计算机学院教师，研究方向为计算机视觉与轻量化模型，指导算法选型与模型优化。
  \item \hlgreen{行业顾问}：通过学校产学研合作渠道对接的基层农技推广人员，为县域推广与统防统治项目运作提供实务指导。
\end{itemize}

\section{组织架构与人员规划}

\begin{itemize}[leftmargin=1.6em]
  \item \hlgreen{初创期（第 1 年）}：5--8 人学生团队，按「算法＋产品＋农技内容＋推广」四条线分工，扁平化运作，研发场地依托学校创新创业平台。
  \item \hlgreen{成长期（第 2--3 年）}：扩充至 10--15 人，吸纳低年级学生与实习人员参与数据标注、渠道拓展与客户维护；团队毕业后视项目进展决定是否全职投入。
\end{itemize}

\section{团队优势与风险说明}

团队的核心优势在于\hlgold{「专业互补、成本可控、院校资源可依托」}：农技内容由植保专业成员负责并接受指导教师把关，保证防治方案的专业性；技术开发由计算机与软件专业成员承担，学校实验室可提供算力与场地支持，使项目在极低的资金投入下即可完成产品原型与试点验证。

同时团队也清楚认识到自身的局限：\hlgreen{成员均为在校学生，缺乏完整的商业运营与规模化渠道管理经验，资金实力有限}。因此本项目在商业规划上刻意采取保守的「小规模试点—验证—再扩张」路径，前期聚焦 1--2 个县域把模型跑通、把模式验证清楚，并依托学校产学研合作与指导教师资源弥补行业经验短板，待模式验证后再引入外部资金与全职人才加速扩张。
